\documentclass[11pt,a4paper]{article}
\usepackage[utf8]{inputenc}
\usepackage[T1]{fontenc}
\usepackage[czech]{babel}
\usepackage[margin=2.0cm,top=2.0cm,bottom=2.2cm]{geometry}
\usepackage{amsmath,amssymb}
\usepackage{parskip}
\usepackage{enumitem}
\setlist{nosep,leftmargin=1.4em}
\usepackage{booktabs}
\usepackage{array}
\usepackage{xcolor}
\definecolor{linkblue}{RGB}{20,60,150}
\definecolor{blokcol}{RGB}{20,70,140}
\definecolor{mostcol}{RGB}{20,110,80}
\definecolor{defbg}{RGB}{240,244,252}
\definecolor{defframe}{RGB}{100,130,190}
\definecolor{markcol}{RGB}{176,84,0}
\definecolor{markbg}{RGB}{253,246,236}
\usepackage[most]{tcolorbox}
\newtcolorbox{nitbox}[1][]{breakable,colback=defbg,colframe=defframe,boxrule=0.6pt,arc=1.5mm,
  left=2mm,right=2mm,top=1mm,bottom=1mm,title={#1},fonttitle=\bfseries}
\newtcolorbox{poznbox}[1][]{breakable,colback=markbg,colframe=markcol,boxrule=0.6pt,arc=1.5mm,
  left=2mm,right=2mm,top=1mm,bottom=1mm,title={#1},fonttitle=\bfseries}
\usepackage[colorlinks=true,linkcolor=linkblue,urlcolor=linkblue]{hyperref}

\setlength{\parskip}{0.35em}

% --- sazba osnov ---------------------------------------------------
% \blok{číslo}{název bloku} + \tema{jména témat oddělená \tb}
\newcommand{\tb}{\ \textcolor{defframe}{\textbullet}\ }
\newcommand{\blok}[2]{\par\addvspace{2.4mm}\noindent
  \textcolor{blokcol}{\textbf{#1}}\ \ \textbf{#2}\par\nobreak}
\newcommand{\tema}[1]{\par\nobreak\noindent #1\par}
% \most{...} = jednou větou, čím se blok přiváže na následující
\newcommand{\most}[1]{\par\nobreak\noindent
  {\small\itshape\textcolor{mostcol}{$\hookrightarrow$\ #1}}\par}
% \apl{metoda}{oblast využití} + 1--2 věty inline (jen u~otázky o~praktickém využití)
\newcommand{\apl}[2]{\par\addvspace{1.5mm}\noindent
  \textbf{#1}\ \textcolor{defframe}{$\rightarrow$}\ \textbf{#2}.\ \ignorespaces}
% \lbl{...} = malý oranžový popisek, text pokračuje inline
\newcommand{\lbl}[1]{\par\addvspace{1.8mm}\noindent
  \textcolor{markcol}{\textbf{\small #1}}\ \ignorespaces}
% oficiální znění věty okruhu pod nadpisem otázky
\newcommand{\zneni}[1]{\par\nobreak\vspace{-1.2mm}\noindent
  {\small\itshape\uv{#1}}\par\addvspace{1.2mm}}

% oddělovač okruhů: #1 = číslo zdrojového PDF, #2 = název okruhu, #3 = podtitul
\newcommand{\okruh}[3]{%
  \clearpage
  \setcounter{section}{0}%
  \phantomsection
  \addcontentsline{toc}{part}{#1\quad #2}%
  \begingroup\centering
    {\Huge\bfseries #1\quad #2\par}\vspace{1.5mm}%
    {\small\itshape #3\par}%
  \endgroup
  \vspace{4mm}}

\makeatletter
\renewcommand\section{\@startsection{section}{1}{\z@}%
  {-2.6ex \@plus -1ex \@minus -.2ex}{1.0ex \@plus.2ex}{\normalfont\Large\bfseries}}
\renewcommand\subsection{\@startsection{subsection}{2}{\z@}%
  {-2.0ex \@plus -1ex \@minus -.2ex}{0.6ex \@plus.2ex}{\normalfont\large\bfseries}}
\makeatother

\setcounter{tocdepth}{1}
\setcounter{secnumdepth}{2}

\title{\textbf{Strategie výkladu ke~SZZ}\\[2mm]
\large Osnovy patnáctiminutových odpovědí na velké otázky}
\author{SZZ Umělá inteligence, MFF UK}
\date{srpen 2026}

\begin{document}
\maketitle

\begin{center}
\begin{minipage}{0.93\textwidth}
\small
Jedna \textbf{velká otázka} $=$ jedna \texttt{\textbackslash section} v~učebních textech
\texttt{01}--\texttt{03}. Na každou z~nich se má dát mluvit \textbf{patnáct minut} souvisle
a~s~logickou návazností. Tento text neobsahuje látku --- ta je v~příslušném textu
a~v~\texttt{04-vzorce} --- ale \textbf{trasu}: v~jakém pořadí témata otevřít, čím jeden blok
navázat na druhý, co si nechat jako rezervu a~co naopak obětovat, aby se odpověď do patnácti
minut vešla. U~každého bloku jsou jen \emph{jména témat}; co k~nim říct, je v~odpovídající
podkapitole zdrojového textu.

Text pokrývá všechny tři okruhy. Nejvíc práce je odvedeno u~otázek, jejichž \emph{znění
nejmenuje, co má odpověď obsahovat} (např. \uv{Problém hledání cesty} nebo \uv{Rojové
optimalizační algoritmy}) --- tam osnovu určuje tenhle text. U~otázek s~vyjmenovaným
obsahem stačí držet pořadí a~přechody.
\end{minipage}
\end{center}

\vspace{2mm}

%=====================================================================
\section*{Jak se staví patnáctiminutová odpověď}
\addcontentsline{toc}{section}{Jak se staví patnáctiminutová odpověď}

\begin{enumerate}[label=\arabic*.,leftmargin=1.8em]
\item \textbf{Nejdřív anons, pak výklad.} První dvě věty: co téma je a~jak ho projdeš ---
      nahlas jmenovat bloky osnovy. Komise pak ví, že máš plán, a~nepřerušuje po první
      minutě. Bez anonsu vypadá i~dobrá odpověď jako přeskakování.
\item \textbf{Nikdy nezačínat výčtem metod.} Vždy: \emph{formulace úlohy} $\to$
      \emph{co je vzor a~co výstup} $\to$ \emph{jak se hledá} $\to$ \emph{jak se pozná, že
      je dobrý}. Tenhle skelet je univerzální záchrana, když se ztratíš.
\item \textbf{Návaznost $=$ slabina předchozího je motivace pro následující.} Nespojovat
      témata slovem \uv{dále}, ale kauzálně: \emph{Apriori generuje kandidáty} $\Rightarrow$
      \emph{proto FP-Growth}; \emph{accuracy selže na nevyvážených datech} $\Rightarrow$
      \emph{proto ROC}. Přechody říkat nahlas --- právě ony odlišují výklad od odříkání.
\item \textbf{Hloubka u~jednoho reprezentanta, šíře u~ostatních.} Z~každé rodiny jeden
      algoritmus doopravdy (mechanismus, cena, kde selže), ostatní jednou větou \uv{tenhle
      řeší to a~to}. Snaha probrat všechno stejně hluboko je nejčastější důvod, proč
      odpověď nedojde do konce.
\item \textbf{Jeden příklad nebo obrázek na otázku.} Malá tabulka transakcí, box plot,
      matice záměn, ROC, tři uzly a~most --- kreslit na tabuli. Dokazuje, že látku umíš
      použít, a~dává čas na rozmyšlení dalšího bloku.
\item \textbf{Držet si rezervu.} Dva až tři pojmy schválně nezmínit a~mít je připravené na
      doplňující dotaz. Doplňující otázka pak přijde na půdu, kterou máš zvládnutou.
\item \textbf{Uzávěr.} Na konci rekapitulovat osnovu jednou větou a~zasadit téma do většího
      obrázku (typicky: kam patří v~procesu KDD, na co navazuje v~jiné otázce). Odpověď,
      která skončí uprostřed technického detailu, působí nedokončeně, i~když byla obsáhlá.
\end{enumerate}

\begin{poznbox}[Když je otázka větší než patnáct minut]
\small
U~nejobsáhlejších otázek (\emph{Metody pro dobývání znalostí}, \emph{Analýza sociálních
sítí}) je úplné pokrytí nemožné. Správný postup je to \textbf{přiznat v~anonsu} a~nabídnout
volbu: \uv{Osnova má pět bloků; do hloubky půjdu u~$\langle$X$\rangle$
a~$\langle$Y$\rangle$, zbytek zmíním --- pokud chcete jinak, přesměrujte mě.} Komise tím
dostane kontrolu nad časem a~ty záruku, že hloubku předvedeš na látce, kterou umíš nejlépe.
\end{poznbox}

%=====================================================================
\okruh{01}{Multiagentní systémy}{Sedm velkých otázek podle oficiálního znění okruhu
  --- kostry odpovědí}

\section*{Přehled: kostry odpovědí}
\addcontentsline{toc}{section}{Přehled: kostry odpovědí (okruh 01)}

\begin{center}
\small
\begin{tabular}{@{}>{\raggedright\arraybackslash}p{3.0cm}>{\raggedright\arraybackslash}p{13.4cm}@{}}
\toprule
\textbf{Otázka} & \textbf{Kostra (bloky za sebou)} \\
\midrule
1 Architektura & co je agent $\to$ prostředí $\to$ abstraktní architektura $\to$ utilita
  a~specifikace úlohy $\to$ mechanismus výběru akcí \\
2 Řízení agentů & deliberativní pól $\to$ BDI a~plánování $\to$ symbolické reaktivní
  $\to$ konekcionistické reaktivní $\to$ omezení reaktivních $\to$ hybridní \\
3 Hledání cesty & formulace $\to$ A* $\to$ dynamické prostředí $\to$ MAPF $\to$ praxe
  a~vazba na MAS \\
4 Komunikace & proč je jiná $\to$ ontologie a~znalosti $\to$ řečové akty $\to$ KQML
  a~FIPA-ACL $\to$ protokoly \\
5 Kooperace a~hry & koherence a~CDPS $\to$ mechanismy kooperace $\to$ sobecký agent a~hry
  $\to$ aukce $\to$ návrh mechanismů \\
6 Učení & mapa učení $\to$ MDP $\to$ učení bez modelu $\to$ aproximace a~politika $\to$
  multiagentní RL \\
7 Metodologie & proč ne OO $\to$ metodologie návrhu $\to$ FIPA a~platforma $\to$ jazyky
  a~prostředí \\
\bottomrule
\end{tabular}
\end{center}

%---------------------------------------------------------------------
\section{Architektura autonomního agenta a~mechanismus výběru akcí}
\zneni{Architektura autonomního agenta; mechanismus výběru akcí.}

\begin{nitbox}[Červená nit]
Znění vypadá jako dvě otázky, ale je to jedna linka: nejdřív \emph{co je agent a~v~jakém
prostředí žije} (formální rámec), pak \emph{jak z~vjemu vznikne akce}. Mechanismus výběru
akcí je přesně to místo, ve kterém se pak liší všechny architektury z~otázky~2 --- tím se
odpověď zakončí a~zároveň připraví půda pro doplňující dotazy.
\end{nitbox}

\blok{1}{Co je agent}
\tema{agent $=$ autonomní systém zasazený do prostředí, který v~něm vnímá a~jedná, aby
splnil delegované cíle
\tb MAS $=$ interagující agenti \emph{bez nutně sdíleného cíle}
\tb tři vlastnosti inteligentního agenta: reaktivita, proaktivita, sociálnost --- a~proč je
těžké je vyvážit
\tb agent vs.\ objekt (vlastní vlákno řízení, autonomie nad přijetím požadavku)
\tb intencionální systém a~tři Dennettovy postoje; intencionální postoj jako legitimní
abstrakce}
\most{Chování agenta má smysl jen vůči prostředí --- druhá polovina rámce je tedy prostředí.}

\blok{2}{Prostředí a~jeho vlastnosti}
\tema{čtyři dichotomie: pozorovatelnost, statické vs.\ dynamické, determinismus, diskrétnost
\tb reálná prostředí jsou v~nejtěžší kombinaci
\tb co která vlastnost zakazuje --- např. dynamika ruší platnost jednou spočítaného plánu}
\most{Aby se dalo o~agentovi něco dokázat, potřebuju formalismus.}

\blok{3}{Abstraktní architektura agenta}
\tema{prostředí $\mathit{Env}=\langle E,e_0,\tau\rangle$, běh $r=e_0a_0e_1\dots$
\tb $\tau: R^A \to 2^E$ --- odtud závislost na historii i~nedeterminismus
\tb agent $Ag: R^E \to A$, systém $\langle Ag,\mathit{Env}\rangle$
\tb funkční ekvivalence agentů (prostá vs.\ vzhledem k~prostředí)
\tb čistě reaktivní agent $Ag: E \to A$ (slabší) vs.\ agent s~vnitřním stavem
\tb role funkcí: $\mathit{see}$ vjem, $\mathit{next}$ \emph{aktualizace stavu},
$\mathit{action}$ výběr akce}
\most{Formalismus říká, \emph{jak} agent jedná; nic zatím neříká, \emph{co} má dělat.}

\blok{4}{Jak agentovi říct, co má dělat}
\tema{utilita nad stavy vs.\ nad běhy --- a~proč stavy nestačí (ignorují průběh)
\tb proč se neohodnocují rovnou akce (prozradilo by řešení)
\tb optimální agent maximalizuje \emph{očekávanou} utilitu; definice je nekonstruktivní
\tb Tileworld a~$u(r)=N_f/N_a$
\tb predikátová specifikace úlohy $\langle \mathit{Env},\Psi\rangle$; pesimista, optimista,
realista
\tb úlohy dosažení a~udržení}
\most{Zbývá poslední díl: jak se ta akce v~každém kroku konkrétně vybere.}

\blok{5}{Mechanismus výběru akcí (jádro otázky)}
\tema{co to je: mapování vnitřní stav $+$ vjem $\to$ akce
\tb rodiny mechanismů a~jejich zástupci: logická dedukce (deduktivní agent) \tb praktické
uvažování a~plánování (BDI, STRIPS, PRS) \tb pevná priorita chování (subsumpce) \tb úroveň
aktivace v~síti modulů (Maes) \tb mediační funkce ve vrstvách (hybridní) \tb soutěž koalic
(IDA) \tb naučená politika (RL)
\tb osy srovnání: kvalita rozhodnutí vs.\ reakční doba, potřeba modelu světa, návrhová
složitost}

\lbl{Rezerva:} rozdíl prosté a~prostředím podmíněné funkční ekvivalence \tb Tileworld jako
testbed \tb agent vs.\ objekt vs.\ expertní systém.

%---------------------------------------------------------------------
\section{Metody pro řízení agentů}
\zneni{Metody pro řízení agentů; symbolické a~konekcionistické reaktivní plánování,
hybridní přístupy.}

\begin{nitbox}[Červená nit]
Znění jmenuje tři body, ale dávají smysl jen na \textbf{jedné ose od deliberativního pólu
k~reaktivnímu a~zpět} --- proto je pořadí dané a~netřeba nad ním váhat. Ke~každé
architektuře říct stejnou trojici: \emph{jak vybírá akci --- co umí --- kde selže}; selhání
je pokaždé motivací té následující.
\end{nitbox}

\blok{1}{Symbolický (deliberativní) pól}
\tema{klasická UI: symbolická reprezentace $+$ syntaktická manipulace, specifikace $=$
program
\tb deduktivní agent: $\mathit{DB} \vdash_\rho \mathit{Do}(a)$, dvoufázový výběr akce
(dokazatelná $\to$ konzistentní $\to$ nic)
\tb problém transdukce
\tb problém vypočitatelné racionality --- odvozená akce platí pro svět, který už není}
\most{Dokazování je moc pomalé; praktické uvažování ho nahradí strukturou.}

\blok{2}{BDI a~plánování}
\tema{praktické uvažování: deliberace $+$ uvažování o~prostředcích
\tb beliefs (subjektivní a~omylné), desires (mohou být nekonzistentní), intentions (závazky)
\tb proč záměry přetrvávají a~\emph{omezují další deliberaci}
\tb funkce $\mathit{brf}$, $\mathit{options}$, $\mathit{filter}$
\tb \textbf{STRIPS}: deskriptor akce $\langle P_a,D_a,A_a\rangle$, plánovací problém,
plán přípustný vs.\ korektní, obcházení frame problému
\tb BDI smyčka a~role $\mathit{reconsider}$
\tb odhodlání: blind / single-minded / open-minded; bold vs.\ cautious podle rychlosti
prostředí
\tb \textbf{PRS}: knihovna plánů místo syntézy, zásobník záměrů; potomci AgentSpeak/Jason}
\most{Celá tahle větev stojí na modelu světa --- reaktivní přístup ho odmítne.}

\blok{3}{Symbolické reaktivní plánování}
\tema{východiska: situatedness, embodiment, emergence; \uv{inteligence bez reprezentace}
\tb \textbf{subsumpční architektura} (Brooks): chování \emph{situace $\to$ akce}, běží
paralelně, pevná hierarchie priorit s~inhibicí
\tb Steelsův průzkumník Marsu a~role stigmergie
\tb pozor na směr inhibice v~zápisu $b_1 \prec b_2$}
\most{Priority jsou tu pevné --- druhá varianta je nechá vyplynout z~dynamiky sítě.}

\blok{4}{Konekcionistické reaktivní plánování}
\tema{\textbf{agent network architecture} (Maes): kompetenční moduly s~předpoklady, add
a~delete listem
\tb šíření aktivace: následnické, předchůdcovské a~konfliktní spoje
\tb tři zdroje aktivace: stav světa, cíle, inhibice od chráněných cílů
\tb výběr proveditelného modulu s~nejvyšší aktivací nad prahem; snižování prahu, aby síť
nezamrzla
\tb parametry $\gamma/\phi$ ladí proaktivitu vs.\ reaktivitu}
\most{Oba reaktivní přístupy narazí na stejnou zeď.}

\blok{5}{Omezení reaktivních architektur}
\tema{jen lokální informace, žádný model světa ani výhled
\tb neinženýrský návrh: emergentní chování se špatně navrhuje i~ověřuje
\tb neškáluje s~počtem vrstev \tb neschopnost učení}
\most{Odtud přímo plyne motivace hybridních architektur.}

\blok{6}{Hybridní architektury}
\tema{horizontální vrstvení: paralelní vrstvy $+$ \emph{mediační funkce}, až $m^n$
kombinací, úzké hrdlo i~kritický bod selhání
\tb vertikální: one-pass vs.\ two-pass, $n-1$ rozhraní, křehkost
\tb obecné pravidlo: reaktivní vrstva má přednost (bezpečnost před optimalitou)
\tb TouringMachines (horizontální, tři vrstvy $+$ řídicí pravidla) a~InteRRaP (vertikální
two-pass)
\tb IDA a~global workspace: koalice procesů soutěží o~vstup do \uv{vědomí}
\tb sjednocující osa: čtyři typy arbitráže --- pevná priorita, úroveň aktivace, mediační
funkce, soutěž koalic}

\lbl{Rezerva:} Stanley jako reálné nasazení vrstvené architektury \tb vysavačový svět
u~deduktivního agenta \tb proč se BDI implementuje knihovnou plánů místo plánovače.

%---------------------------------------------------------------------
\section{Problém hledání cesty}
\zneni{Problém hledání cesty.}

\begin{nitbox}[Červená nit]
Jednovětné znění, osnovu si tedy určuješ sám. Nejlepší stavba je \textbf{stupňující se řada
předpokladů, které se postupně ruší}: statické známé prostředí (A*) $\to$ neznámé a~měnící
se (D*~Lite) $\to$ víc agentů najednou (MAPF). Každý stupeň $=$ jeden předpoklad, který
realita poruší, a~jeden algoritmus, který na to odpovídá. MAPF je navíc téma bakalářské
práce --- přirozené místo, kam odpověď natáhnout do libovolné hloubky.
\end{nitbox}

\blok{1}{Formulace a~zařazení}
\tema{úloha jako prohledávání stavového prostoru: graf, ceny hran, přípustné akce
\tb co je \uv{mapa}: mřížka, navigační graf, navmesh; diskretizace spojitého prostoru
\tb vztah k~plánování (STRIPS): hledání cesty je speciální případ s~jednoduchými akcemi
\tb neinformované prohledávání jako pozadí (BFS, Dijkstra) a~proč nestačí}
\most{Klíčem k~efektivitě je heuristika --- tedy odhad zbytku cesty.}

\blok{2}{A* a~jeho vlastnosti}
\tema{$f = g + h$; open a~closed list
\tb \textbf{přípustnost} $h \le h^*$ $\Rightarrow$ optimalita
\tb \textbf{konzistence} (trojúhelníková nerovnost) $\Rightarrow$ každý vrchol expandován
nejvýš jednou (nutné s~closed listem)
\tb speciální případy: $h\equiv 0$ dá Dijkstru, weighted A* je $\varepsilon$-suboptimální
\tb heuristiky na mřížce (Manhattan, oktilní, euklidovská) a~dominance heuristik
\tb paměť a~čas: proč se v~praxi sahá po IDA*, hierarchickém plánování, jump point search}
\most{A* předpokládá, že mapu znám a~nemění se --- v~robotice ani ve hře to neplatí.}

\blok{3}{Dynamické a~neznámé prostředí}
\tema{freespace assumption: co neznám, považuji za volné
\tb inkrementální přeplánování místo počítání od nuly
\tb \textbf{LPA*}: hodnoty $g$ a~$\mathit{rhs}$, lokální konzistence $g=\mathit{rhs}$
\tb \textbf{D*~Lite}: prohledávání od cíle, oprava jen zasažené části
\tb kdy se vyplatí opravovat a~kdy přepočítat od začátku}
\most{Druhý předpoklad, který v~multiagentním systému padne, je \uv{jsem tu sám}.}

\blok{4}{Hledání cest pro více agentů (MAPF)}
\tema{formulace: graf, agenti, startovní a~cílové vrcholy, diskrétní čas, čekání jako akce
\tb typy konfliktů: vrcholový, hranový (swap), následný, cyklický
\tb cílové funkce: sum-of-costs vs.\ makespan
\tb optimální MAPF je NP-těžký
\tb prioritizované plánování / Cooperative A* s~rezervační tabulkou --- rychlé, ale neúplné
\tb \textbf{CBS}: nízká úroveň A* s~omezeními, vysoká úroveň strom omezení větvený podle
nalezeného konfliktu; optimální a~úplný
\tb suboptimální varianty a~velké instance (sklady, roboti)}
\most{Zbývá říct, proč je celá tahle otázka v~okruhu o~agentech.}

\blok{5}{Praxe a~vazba na MAS}
\tema{pohyb je nejčastější akce agenta; plánování cesty je subúloha uvažování o~prostředcích
\tb hry a~robotika: hierarchie, vyhlazování cesty, any-angle (Theta*)
\tb centralizované plánování (MAPF) vs.\ decentralizované vyhýbání (ORCA) --- a~co to říká
o~koordinaci obecně}

\lbl{Rezerva:} důkaz optimality A* při přípustné heuristice \tb proč konzistence implikuje
přípustnost \tb jak se z~konfliktu stane omezení v~CBS.

%---------------------------------------------------------------------
\section{Komunikace a~znalosti v~MAS}
\zneni{Komunikace a~znalosti v~multiagentních systémech, ontologie, řečové akty, FIPA-ACL,
protokoly.}

\begin{nitbox}[Červená nit]
Čtyři jmenované věci na sebe navazují jako \textbf{vrstvy jednoho zásobníku}: čím se mluví
o~světě (ontologie) $\to$ co promluva znamená (řečové akty) $\to$ jak vypadá zpráva (ACL)
$\to$ jak se zprávy skládají do konverzací (protokoly). Odpověď $=$ průchod zásobníkem
zdola nahoru, uvozený tím, proč je komunikace agentů vůbec zvláštní téma.
\end{nitbox}

\blok{1}{Proč je komunikace agentů jiná}
\tema{agent je autonomní: zpráva je \emph{akce}, ne volání metody --- příjemce není povinen
vyhovět
\tb z~toho plyne, co musí zpráva nést: záměr, jazyk obsahu, ontologii, kontext konverzace}
\most{Nejnižší vrstva je sdílený slovník --- bez něj nelze obsah zprávy interpretovat.}

\blok{2}{Ontologie a~znalosti}
\tema{ontologie $=$ explicitní formalizovaný popis části reality (Gruber)
\tb třídy, instance, vlastnosti, relace (zejména tranzitivní \emph{is-a}), fakta; ontologie
$+$ fakta $=$ znalostní báze
\tb škála formalizace od slovníku po libovolná logická omezení; aplikační, doménové a~horní
ontologie
\tb RDF (trojice subjekt--predikát--objekt), RDFS (třídy, subClassOf, domain, range)
\tb \textbf{OWL}: deskripční logika, rozhodnutelný reasoning, T-Box vs.\ A-Box, open world
assumption, profily
\tb KIF jako jazyk obsahu
\tb epistemická logika: $K_i\varphi$ a~možné světy; znalost S5 (axiom~T) vs.\ přesvědčení
KD45; společná znalost a~koordinovaný útok}
\most{Slovník mám; teď co vlastně promluva \emph{dělá}.}

\blok{3}{Teorie řečových aktů}
\tema{Austin: lokuce / ilokuce / perlokuce
\tb Searle: podmínky zdařilosti a~pět kategorií (asertiva, direktiva, komisiva, expresiva,
deklarace)
\tb performativní sloveso $+$ propoziční obsah
\tb Cohen \& Perrault: sémantika ve stylu STRIPS (předpoklady a~efekty s~modálními
operátory); $\mathit{Request}$ a~$\mathit{Inform}$ --- efektem Informu je, že
\emph{příjemce věří, že odesílatel věří} $\varphi$}
\most{To je teorie; standard z~ní udělá konkrétní formát zprávy.}

\blok{4}{KQML a~FIPA-ACL}
\tema{KQML $=$ obálka, KIF $=$ obsah; kritika (chybí komisiva, nejasná sémantika)
\tb \textbf{FIPA-ACL}: 13 polí, povinné je jen \texttt{performative}
\tb sender/receiver (AID), content, language, encoding, ontology, protocol,
conversation-id, reply-with, in-reply-to, reply-by
\tb 22 performativ; sémantika přes feasibility precondition a~rational effect v~jazyce SL
\tb problém neverifikovatelnosti --- vnitřní stav odesílatele nelze ověřit}
\most{Jedna zpráva nestačí; dialog má předepsaný tvar.}

\blok{5}{Interakční protokoly}
\tema{protokol jako předepsaná posloupnost performativ; role zahajovatele a~účastníka
\tb FIPA-Request: \texttt{request} $\to$ \texttt{refuse} $|$ \texttt{agree} $\to$
\texttt{failure} $|$ \texttt{inform-done} $|$ \texttt{inform-result}, kdykoli
\texttt{not-understood}
\tb FIPA-Query a~Subscribe jako tentýž skelet
\tb Contract Net a~jeho iterovaná varianta \tb brokering \tb aukční protokoly
\tb kde protokol \emph{nestačí}: nebrání lhaní --- to už je návrh mechanismů (ot.~5)}

\lbl{Rezerva:} proč má ACL zpráva pole \texttt{ontology} \tb axiomy T, D, 4, 5 a~co který
znamená \tb proč společná znalost nevznikne nespolehlivou komunikací.

%---------------------------------------------------------------------
\section{Distribuované řešení problémů, kooperace, teorie her a~aukce}
\zneni{Distribuované řešení problémů, kooperace, Nashova ekvilibria, Paretova efektivita,
alokace zdrojů, aukce.}

\begin{nitbox}[Červená nit]
Otázka má dvě poloviny s~\textbf{opačným předpokladem}: benevolentní agenti se sdíleným
cílem (distribuované řešení problémů) a~sobečtí agenti s~vlastními cíli (teorie her, aukce).
Pointa, kterou má odpověď doručit: jakmile padne benevolence, přestává stačit protokol
a~začíná \emph{návrh mechanismu}. Tuhle větu říct hned v~anonsu --- drží pak celou odpověď.
\end{nitbox}

\blok{1}{Koherence, koordinace a~CDPS}
\tema{koherence $=$ kvalita chování celku (cíl) vs.\ koordinace $=$ správa vzájemných
závislostí (prostředek)
\tb CDPS: benevolence $+$ sdílený cíl; odlišit od paralelního řešení i~od MAS se sobeckými
agenty
\tb tři fáze: dekompozice $\to$ řešení podproblémů $\to$ syntéza
\tb sdílení úloh vs.\ sdílení výsledků (proaktivní a~reaktivní) a~jeho čtyři přínosy}
\most{Nejznámější protokol pro sdílení úloh je zároveň nejcitovanější v~celém okruhu.}

\blok{2}{Mechanismy kooperace}
\tema{\textbf{Contract Net}: recognition $\to$ announcement $\to$ bidding $\to$ awarding
a~expediting; manažer a~dodavatel, hierarchické subkontrakty
\tb co CNET \emph{negarantuje}: globální optimalitu ani příchod nabídky; předpokládá
benevolenci, nebrání lhaní o~schopnostech
\tb systémy s~tabulí: hierarchie abstrakce, experti, arbitr; volná vazba, úzké hrdlo
\tb DisCSP a~DCOP: proměnné a~omezení rozdělené mezi agenty, nikdo nevidí celý problém;
ABT (\texttt{ok?} / \texttt{nogood}), DCOP maximalizuje sociální blahobyt}
\most{Všechno výše padá, jakmile agent sleduje vlastní zájem --- a~to je druhá polovina
otázky.}

\blok{3}{Sobecký agent, hry a~ekvilibria}
\tema{sobecký agent maximalizuje \emph{svou} očekávanou utilitu za předpokladu, že ostatní
dělají totéž
\tb hra v~normální formě; ryzí vs.\ smíšené strategie; dominance a~iterované odstraňování
dominovaných strategií
\tb \textbf{Nashovo ekvilibrium} $=$ vzájemně nejlepší odpovědi; Nashova věta (existence ve
smíšených strategiích); hledání je PPAD-úplné
\tb \textbf{Paretova optimalita} vs.\ sociální blahobyt; vztahy: profil dominantních
strategií je vždy NE, NE nemusí být Pareto-optimální
\tb vězňovo dilema jako ilustrace všech tří pojmů; tragedy of the commons; iterovaná verze,
stín budoucnosti, TFT
\tb klasické hry: koordinační, anti-koordinační (chicken), s~nulovým součtem
\tb umět dopočítat smíšené NE ve hře $2\times2$ (volím tak, aby byl \emph{protihráč}
indiferentní)}
\most{Když se agenti neshodnou, kdo co dostane, potřebují pravidla --- a~znění okruhu jmenuje
ta tržní.}

\blok{4}{Aukce jako alokace zdrojů}
\tema{aukce $=$ mechanismus alokace vzácného zdroje; efektivní aukce dá zdroj tomu, kdo ho
oceňuje nejvíc
\tb dimenze: private vs.\ common value, open cry vs.\ sealed bid, first vs.\ second price,
počet kol
\tb čtyři základní typy a~jejich strategie; dominantní strategie má anglická a~Vickreyova,
holandská a~FPSB ji nemají
\tb důkaz pravdomluvnosti u~Vickreyho: nabídka rozhoduje jen o~tom, \emph{zda} vyhraješ,
ne kolik zaplatíš
\tb ekvivalence: anglická $\leftrightarrow$ Vickrey, holandská $\leftrightarrow$ FPSB
\tb věta o~ekvivalenci výnosů a~kdy neplatí (averze k~riziku, common values, prokletí
vítěze, koluze)
\tb kombinatorické aukce a~NP-těžký problém určení vítěze}
\most{Poslední krok obrací směr úvahy: nenavrhovat strategii, ale pravidla hry.}

\blok{5}{Návrh mechanismů a~vyjednávání}
\tema{mechanism design jako obrácená teorie her; princip odhalení
\tb \textbf{VCG}: alokace maximalizující sociální blahobyt $+$ platba rovná
\emph{externalitě}; pravdomluvnost jako dominantní strategie, efektivita, individuální
racionalita --- ale není budget-balanced
\tb vyjednávání: monotónní protokol ústupků $+$ Zeuthenova strategie (ustoupí ten s~nižším
rizikem), vazba na Nashův bargaining product}

%---------------------------------------------------------------------
\section{Metody pro učení agentů a~zpětnovazební učení}
\zneni{Metody pro učení agentů; zpětnovazební učení.}

\begin{nitbox}[Červená nit]
Znění jmenuje dvě věci a~druhá je podmnožinou první --- tak to i~postavit: nejdřív
\textbf{mapa toho, co všechno \uv{učení agenta} znamená}, pak \textbf{jedna větev (RL) do
hloubky}, protože je pro agenta nejpřirozenější (dostává odměnu, ne správnou odpověď).
Konec patří tomu, co se rozbije, jakmile je agentů víc --- tím se otázka napojí na okruh
teorie her.
\end{nitbox}

\blok{1}{Mapa učení agentů}
\tema{proč se agent učí: prostředí není známé předem a~návrhář nezná optimální politiku
\tb s~učitelem / bez učitele / zpětnovazební \tb imitační učení a~inverzní RL
\tb evoluční metody (LCS, neuroevoluce) --- vazba na okruh~02
\tb co se agent učí: model prostředí, model protivníka, politiku
\tb individuální vs.\ sociální učení}
\most{RL je pro agenta ta pravá větev: má jen skalární odměnu a~musí sám objevit, co ji
přináší.}

\blok{2}{Formalismus: MDP}
\tema{$\langle S,A,T,R,\gamma\rangle$ a~markovská vlastnost
\tb politika, hodnotové funkce $V^\pi$ a~$Q^\pi$, role diskontování
\tb Bellmanova rovnice optimality
\tb value iteration (kontrakce, asymptotická konvergence) vs.\ policy iteration (evaluace
$+$ zlepšení, konečné ukončení)}
\most{Obojí ale potřebuje model prostředí --- a~ten agent typicky nemá.}

\blok{3}{Učení bez modelu}
\tema{model-based vs.\ model-free
\tb \textbf{Q-learning}: aktualizační pravidlo, off-policy, tabulkový; umět projít jeden
krok s~čísly
\tb \textbf{SARSA}: on-policy a~proč je opatrnější (cliff walking)
\tb dilema průzkum/využití: $\varepsilon$-greedy vs.\ \textbf{UCB} (bonus za nejistotu,
\uv{optimismus tváří v~tvář nejistotě}, bonus sám vyhasíná)}
\most{Tabulka přestane stačit, jakmile je stavů příliš mnoho.}

\blok{4}{Aproximace a~metody založené na politice}
\tema{aproximace hodnotové funkce; DQN (replay buffer, cílová síť) a~proč jsou obě potřeba
\tb REINFORCE: gradient logaritmu politiky krát výnos
\tb actor--critic: herec $=$ politika, kritik $=$ hodnotová funkce, výhoda $A = Q - V$
\tb kdy politika místo hodnoty: spojité akce, stochastické politiky}
\most{Všechno dosud předpokládalo, že se prostředí nemění samo od sebe --- v~MAS se mění.}

\blok{5}{Multiagentní zpětnovazební učení}
\tema{markovská hra: sdružený akční prostor, individuální odměny
\tb místo optima se hledá řešení: nejlepší odpověď, Nashovo ekvilibrium, Pareto (vazba na
ot.~5)
\tb Minimax-Q (nulový součet, LP v~každé aktualizaci), Nash-Q
\tb čtyři problémy MARL: nestacionarita, exponenciální sdružený prostor, částečná
pozorovatelnost, přiřazení zásluh
\tb CTDE jako standardní odpověď \tb self-play a~AlphaZero}

\lbl{Rezerva:} proč value iteration konverguje (kontrakce) \tb rozdíl mezi odměnou
a~utilitou běhu (vazba na ot.~1) \tb v~čem přesně je Q-learning off-policy.

%---------------------------------------------------------------------
\section{Metodologie návrhu, jazyky a~prostředí MAS}
\zneni{Metodologie návrhu, jazyky a~prostředí multiagentních systémů.}

\begin{nitbox}[Červená nit]
Tři slova ve znění $=$ tři vrstvy jedné otázky: \textbf{jak systém navrhnout} (metodologie)
$\to$ \textbf{podle čeho ho postavit} (standard FIPA) $\to$ \textbf{v~čem ho napsat}
(jazyky a~platformy). Je to nejfaktografičtější otázka okruhu --- nedá se v~ní improvizovat,
zato se dá spolehlivě naučit. Začít tím, \emph{proč nestačí objektově orientovaný návrh}.
\end{nitbox}

\blok{1}{Proč nestačí objektově orientované metodologie}
\tema{autonomie: agent může požadavek odmítnout, objekt ne
\tb vlastní cíle a~organizační (nikoli strukturální) vztahy
\tb důsledek: analýza pracuje s~rolemi a~interakcemi, ne s~třídami a~metodami}
\most{Nejznámější metodologie staví přesně na pojmu role.}

\blok{2}{Metodologie návrhu}
\tema{\textbf{Gaia}: analýza $=$ model rolí (permissions, responsibilities s~liveness
a~safety, activities, protocols) $+$ model interakcí; návrh $=$ model agentů, služeb
a~známostí; předpoklady --- kooperativní agenti a~statická organizace
\tb rozšíření ROADMAP a~Gaia v.2, která tyhle předpoklady uvolňují}
\most{Návrh je hotový; aby spolu systémy mluvily, potřebují standard.}

\blok{3}{FIPA a~architektura platformy}
\tema{co FIPA standardizuje: ACL, protokoly, architekturu platformy
\tb \textbf{AMS} --- bílé stránky, přiděluje AID, řídí životní cyklus (initiated, active,
waiting, suspended, transit)
\tb \textbf{DF} --- žluté stránky: registrace a~vyhledávání \emph{služeb}
\tb \textbf{MTS/ACC} --- doručování ACL zpráv v~rámci platformy i~mezi platformami
\tb rozdíl standard vs.\ implementace (FIPA vs.\ JADE) --- říct explicitně}
\most{Standard popisuje rozhraní; kód se píše v~konkrétním prostředí.}

\blok{4}{Jazyky a~prostředí}
\tema{\textbf{JADE}: Java, plně FIPA; třída \texttt{Agent} se \texttt{setup()}
a~\texttt{takeDown()}; chování jako objekty \texttt{Behaviour} (OneShot, Cyclic, Ticker,
Sequential, FSM) s~\texttt{action()} a~\texttt{done()}
\tb jedno vlákno na agenta a~v~něm \emph{kooperativní nepreemptivní scheduler}
$\Rightarrow$ blokující chování zastaví celého agenta, odtud nutnost \texttt{block()}
\tb \texttt{ACLMessage} a~\texttt{MessageTemplate}, hotové protokoly
(\texttt{ContractNetInitiator}), containers
a~mobilita, nástroje RMA a~Sniffer; registrace služby u~DF v~\texttt{setup()}
\tb \textbf{Jason/AgentSpeak}: BDI jazyk, plán \texttt{spouštěč : kontext <- tělo};
\texttt{!g} vs.\ \texttt{?g}, \texttt{tell}/\texttt{achieve} jako řečové akty, anotace
\texttt{[source(...)]}, \texttt{not} vs.\ silná negace; prostředí v~Javě, projekt
\texttt{.mas2j}
\tb \textbf{NetLogo}: agentní simulace (turtles, patches, links) --- jiný účel než JADE
\tb \textbf{SPADE}: Python $+$ XMPP
\tb jak vybrat: JADE pro interoperabilitu, Jason pro BDI abstrakci, NetLogo pro simulaci}

\lbl{Rezerva:} co konkrétně se píše do \texttt{setup()} \tb selekční funkce cyklu
a~\texttt{AgArch} v~Jasonu \tb proč Gaia nepočítá se sobeckými agenty.

%=====================================================================
\okruh{02}{Přírodou inspirované počítání}{Šest velkých otázek podle oficiálního znění
  okruhu --- kostry odpovědí}

\section*{Přehled: kostry odpovědí}
\addcontentsline{toc}{section}{Přehled: kostry odpovědí (okruh 02)}

\begin{center}
\small
\begin{tabular}{@{}>{\raggedright\arraybackslash}p{3.0cm}>{\raggedright\arraybackslash}p{13.4cm}@{}}
\toprule
\textbf{Otázka} & \textbf{Kostra (bloky za sebou)} \\
\midrule
1 GA, GP, EP & společné schéma EA $\to$ genetický algoritmus $\to$ reprezentace a~operátory
  $\to$ genetické programování $\to$ evoluční programování \\
2 Teorie schémat & schéma a~jeho míry $\to$ věta o~schématech $\to$ stavební bloky $\to$
  kritika $\to$ Voseho model $\to$ konečné populace \\
3 ES, DE, koevoluce & evoluční strategie $\to$ diferenciální evoluce $\to$ koevoluce $\to$
  otevřená evoluce \\
4 Rojové algoritmy & princip rojové inteligence $\to$ PSO $\to$ ACO $\to$ další roje $\to$
  srovnání s~EA \\
5 Lokální a~memetické & hill climbing $\to$ simulované žíhání $\to$ tabu a~scatter search
  $\to$ memetické algoritmy $\to$ parametry \\
6 Aplikace & pravidlové systémy a~LCS $\to$ neuroevoluce $\to$ kombinatorické úlohy
  a~omezení $\to$ vícekriteriální optimalizace \\
\bottomrule
\end{tabular}
\end{center}

%---------------------------------------------------------------------
\section{Genetické algoritmy, genetické a~evoluční programování}
\zneni{Genetické algoritmy, genetické a~evoluční programování.}

\begin{nitbox}[Červená nit]
Tři jmenované věci jsou \textbf{tři body na jedné ose: co se vlastně evolvuje} --- řetězec
(GA), program jako strom (GP), automat či chování (EP). Postavit odpověď na \emph{společném
schématu evolučního algoritmu} a~pak ukázat, že se větve liší jen volbou reprezentace
a~operátorů. Šetří to čas a~drží odpověď pohromadě; jinak se otázka rozpadne na tři
nesouvisející výčty.
\end{nitbox}

\blok{1}{Společné schéma evolučního algoritmu}
\tema{EA $=$ populační stochastické prohledávání
\tb cyklus: inicializace $\to$ ohodnocení $\to$ rodičovská selekce $\to$ křížení $\to$
mutace $\to$ ohodnocení $\to$ environmentální selekce
\tb variace tvoří diverzitu, selekce ji odebírá $\Rightarrow$ explorace vs.\ exploatace jako
hlavní návrhová osa
\tb No Free Lunch $\Rightarrow$ doménová specializace reprezentace a~operátorů
\tb historické větve GA / ES / EP / GP a~jejich charakteristické volby}
\most{Nejjednodušší instancí schématu je jednoduchý genetický algoritmus.}

\blok{2}{Genetický algoritmus}
\tema{SGA: binární řetězce $+$ ruleta $+$ jednobodové křížení $+$ bit-flip mutace $+$
generační náhrada; typické hodnoty $p_C$ a~$p_M$
\tb selekce: ruleta (vysoký rozptyl, citlivá na škálování), SUS, turnaj (tlak přes $k$,
invariantní k~transformacím fitness), pořadová, Boltzmannova
\tb elitismus a~monotonie nejlepší fitness; generační vs.\ steady-state model
\tb škálování fitness (lineární, sigma truncation, mocninné, Boltzmannovo) --- lék na
dominanci \uv{supermanů} na začátku a~ztrátu tlaku na konci}
\most{Všechno tohle je nezávislé na tom, co je v~genomu --- a~právě genom rozhoduje
o~úspěchu.}

\blok{3}{Reprezentace a~genetické operátory}
\tema{reprezentace $+$ operátory $=$ fitness krajina
\tb binární kódování reálných čísel, Hammingovy útesy, Grayův kód
\tb křížení: jednobodové (rozbití schématu), dvoubodové (bez okrajového biasu), uniformní
(nejvíc disruptivní, bez pozičního biasu)
\tb mutace je v~GA sekundární: brání nevratné ztrátě alel
\tb reálná reprezentace: gaussovská mutace, aritmetické křížení a~jeho smršťování rozptylu
$\Rightarrow$ BLX-$\alpha$ a~SBX
\tb permutace: PMX, OX, CX, ER --- co který zachovává; mutace swap, insert, inverze (2-opt),
scramble
\tb nominální vs.\ ordinální atributy (nezaujatá mutace vs.\ creep)}
\most{Když se místo hodnot evolvuje program, mění se všechno kromě samotného schématu.}

\blok{4}{Genetické programování}
\tema{stromová reprezentace: terminály a~neterminály, uzavřenost a~dostatečnost
\tb inicializace grow / full / ramped half-and-half
\tb křížení výměnou podstromů, typy mutací (včetně mutace konstant); fitness $=$ spuštění
programu na sadě případů
\tb \textbf{bloat}: definice, tři vysvětlení (replication accuracy, removal bias, crossover
bias), obrana (limity velikosti, parsimony pressure, anti-bloat operátory, vícekriteriální
přístup)
\tb modularita: ADF; typované a~gramatické GP
\tb varianty: lineární GP, grafové a~Cartesian GP, vývojové GP
\tb gramatická evoluce: celočíselný genom, výběr pravidla modulem, derivační strom;
robustní kódování, ale nelokální mapování}
\most{Poslední větev nechává jako genotyp automat a~křížení zahazuje úplně.}

\blok{5}{Evoluční programování}
\tema{konečné automaty, genotyp $=$ fenotyp, \emph{jen mutace}
\tb turnajová selekce z~rodičů i~potomků
\tb příklad predikce a~degenerace fitness
\tb moderní EP $=$ reálné vektory se samoadaptací (blízko evolučním strategiím)
\tb headless chicken experiment: kombinuje křížení stavební bloky, nebo je to jen
makromutace?}

%---------------------------------------------------------------------
\section{Teorie schémat a~pravděpodobnostní modely SGA}
\zneni{Teorie schémat, pravděpodobnostní modely jednoduchého genetického algoritmu.}

\begin{nitbox}[Červená nit]
Znění jmenuje \textbf{dvě odpovědi na jednu otázku: proč GA funguje}. První (schémata) je
intuitivní, ale je to jen nerovnost pro jeden krok a~má silnou kritiku; druhá (Voseho model
a~Markovovy řetězce) je exaktní, ale prakticky nepoužitelná. Tuhle dvojici postavit proti
sobě --- ten protiklad \emph{je} celá odpověď a~zároveň nejlepší závěrečná věta.
\end{nitbox}

\blok{1}{Schéma a~jeho míry}
\tema{schéma $=$ slovo nad $\{0,1,*\}$
\tb řád $o(S)$ $=$ počet pevných pozic; definiční délka $d(S)$; fitness $F(S)$ jako průměr
reprezentantů \emph{v~dané populaci}
\tb celkem $3^m$ schémat, jeden jedinec je reprezentantem $2^m$ z~nich}
\most{Otázka zní, jak se počet reprezentantů schématu vyvíjí mezi generacemi.}

\blok{2}{Věta o~schématech}
\tema{odvození ve třech krocích: selekce (exponenciální růst nadprůměrných schémat)
$\to$ křížení (rozbití s~pravděpodobností $d(S)/(m-1)$) $\to$ mutace ($(1-p_M)^{o(S)}$)
\tb výsledná nerovnost a~jak ji číst
\tb co věta říká --- a~co neříká: jen dolní odhad na jeden krok}
\most{Z~věty se ale udělal mnohem silnější závěr o~tom, jak GA hledá.}

\blok{3}{Hypotéza stavebních bloků a~implicitní paralelismus}
\tema{BBH: GA skládá krátká nadprůměrná schémata nízkého řádu
\tb důsledky pro praxi: záleží na kódování (geny, které patří k~sobě, mají být blízko), na
velikosti populace; škodí předčasná konvergence
\tb implicitní paralelismus ($\sim n^3$ užitečně zpracovaných schémat)
\tb analogie s~$k$-rukým banditou: exponenciální alokace pokusů vítězi}
\most{A~právě tady začíná kritika, kterou je potřeba umět vyjmenovat.}

\blok{4}{Kritika teorie schémat}
\tema{nerovnost platí jen pro jeden krok \tb konečné populace
\tb deceptivní úlohy (royal road, trap funkce)
\tb kolaterální konvergence a~velký rozptyl fitness
\tb ignoruje konstruktivní efekty operátorů \tb schémata nejsou nezávislá}
\most{Odtud snaha nahradit nerovnost exaktním modelem.}

\blok{5}{Voseho model (nekonečné populace)}
\tema{řetězce ztotožněné s~čísly; populace $=$ vektor relativních četností v~simplexu
\tb hledá se $\mathcal{G} = \mathcal{M} \circ \mathcal{F}$
\tb selekce $\mathcal{F}$: diagonální matice fitness
\tb míchání $\mathcal{M}$ (křížení $+$ mutace): kvadratický operátor; díky symetrii stačí
\emph{jediná} míchací matice
\tb pevné body: selekce $\to$ populace kopií nejlepšího, míchání $\to$ rovnoměrné rozdělení;
jejich kombinace vede k~přerušovaným rovnováhám}
\most{Model je deterministický, protože předpokládá nekonečnou populaci --- realita je
konečná.}

\blok{6}{Konečné populace a~Markovův řetězec}
\tema{stavy $=$ populace; počet stavů je kombinační číslo, tedy astronomický
\tb matice přechodu je exaktní, ale nepoužitelná
\tb ireducibilita při $p_M>0$; elitistický GA konverguje k~optimu s~pravděpodobností 1
\tb závěrečný protiklad: nekonečné populace $=$ exaktní a~nepraktické, konečné $=$
realistické a~nezvládnutelné; limitní přechod obojí spojuje}

\lbl{Rezerva:} spočítat $o(S)$ a~$d(S)$ na konkrétním schématu \tb proč elitismus zaručuje
konvergenci \tb deceptivní funkce jako protipříklad k~BBH.

%---------------------------------------------------------------------
\section{Evoluční strategie, diferenciální evoluce, koevoluce, otevřená evoluce}
\zneni{Evoluční strategie, diferenciální evoluce, koevoluce, otevřená evoluce.}

\begin{nitbox}[Červená nit]
Čtyři jmenované věci nejsou variace téhož a~vyplatí se to říct hned: první dvě jsou
\textbf{algoritmy pro reálnou optimalizaci} a~liší se tím, \emph{odkud berou měřítko kroku}
(strategické parametry vs.\ rozptyl populace); druhé dvě mění \textbf{samotný pojem fitness}
--- u~koevoluce je kontextová, u~otevřené evoluce žádná. Ta dělicí čára vysvětlí, proč jsou
v~jedné otázce.
\end{nitbox}

\blok{1}{Evoluční strategie}
\tema{reálné vektory; mutace je hlavní operátor, rodičovská selekce náhodná,
environmentální deterministická; jedinec $= [\vec x, \vec\sigma]$
\tb $(\mu,\lambda)$ vs.\ $(\mu+\lambda)$: elitismus, konvergence, únik z~lokálních optim,
samoadaptace
\tb $(1+1)$-ES a~pravidlo $1/5$ --- adaptivní, nikoli samoadaptivní řízení
\tb samoadaptace: \textbf{nejdřív mutovat $\sigma$, pak $x$}; počty strategických parametrů
(izotropní, nekorelované, korelované)
\tb rekombinace: lokální vs.\ globální, diskrétní vs.\ intermediární
\tb \textbf{CMA-ES}: kovarianční matice, evoluční cesty, rank-1 a~rank-$\mu$ update, řízení
délky kroku; invariance vůči monotónním transformacím fitness}
\most{ES si měřítko kroku pamatuje v~parametrech; DE ho čte přímo z~populace.}

\blok{2}{Diferenciální evoluce}
\tema{DE/rand/1/bin: dárce $\vec v = \vec x_a + F(\vec x_b - \vec x_c)$, binomické křížení
s~prahem $C$ a~pojistkou, souboj potomka \emph{s~vlastním rodičem}
\tb geometrická podstata: měřítko i~orientace kroků se přebírají z~rozptylu populace
$\Rightarrow$ adaptace bez strategických parametrů
\tb parametry $F$, $C$, velikost populace
\tb varianty rand/1, best/1, current-to-best; \texttt{bin} vs.\ \texttt{exp} křížení
\tb umět předvést jeden krok s~čísly (mutace $\to$ křížení $\to$ selekce)}
\most{Obě metody předpokládají pevnou fitness funkci --- koevoluce ji nechá záviset na
populaci.}

\blok{3}{Koevoluce}
\tema{kontextová fitness $\Rightarrow$ dynamická krajina
\tb \emph{kompetitivní}: predátor--kořist, host--parazit, Hillisovy řadicí sítě
\tb \emph{kooperativní}: dekompozice na komponenty, CCGA, moduly $+$ blueprinty
\tb patologie: cyklení (intranzitivita), odpojení, efekt Červené královny, průměrné stabilní
stavy, přespecializace, zapomínání
\tb léky: archivy (hall of fame), sdílená fitness, měření proti pevné sadě
\tb iterované vězňovo dilema jako klasická úloha (TFT a~vlastnosti úspěšné strategie)
\tb udržování diverzity: fitness sharing, crowding, speciace, ostrovní model}
\most{Poslední krok zahodí cíl úplně.}

\blok{4}{Otevřená evoluce}
\tema{evoluce bez pevného cíle; čtyři znaky --- trvalá novost, neomezený růst složitosti,
nové třídy entit, evoluce evolvability
\tb systémy Tierra, Avida, Polyworld, ECHO a~jejich společný osud (dříve či později se
zaseknou)
\tb generativní (nepřímé) reprezentace: L-systémy, celulární automaty, buněčné kódování,
CPPN
\tb \textbf{novelty search}: řídkost v~prostoru chování vůči populaci \emph{a} archivu;
účelová funkce se nepoužívá vůbec
\tb navazující quality-diversity a~MAP-Elites
\tb interaktivní evoluce (fitness $=$ člověk) a~její limity}

\lbl{Rezerva:} proč se u~samoadaptace mutuje nejdřív $\sigma$ \tb dynamická fitness
(diploidie, hypermutace, migrace) \tb proč $(\mu,\lambda)$ lépe uniká z~lokálních optim.

%---------------------------------------------------------------------
\section{Rojové optimalizační algoritmy}
\zneni{Rojové optimalizační algoritmy.}

\begin{nitbox}[Červená nit]
Znění je jedna věta, osnovu si tedy určuješ sám. Nejsilnější stavba: nejdřív \textbf{princip
společný celé rodině} (decentralizace, stigmergie, kladná zpětná vazba s~tlumením,
emergence), pak \textbf{PSO a~ACO jako dvě různá provedení téhož principu} --- jedno ve
spojitém prostoru, druhé na grafu --- a~nakonec \textbf{srovnání s~evolučními algoritmy},
protože ta otázka přijde tak jako tak. Sama látka je krátká, tenhle rámec ji roztáhne na
patnáct minut, aniž by se plácalo prázdno.
\end{nitbox}

\blok{1}{Princip rojové inteligence}
\tema{inspirace: hejno, mravenčí kolonie, včelstvo; jednoduchý jedinec, žádné centrální
řízení
\tb klíčové pojmy: decentralizace, \textbf{stigmergie} (komunikace přes prostředí), kladná
zpětná vazba $+$ tlumení, emergence
\tb proč to funguje: kladná zpětná vazba zesiluje dobré řešení, odpařování a~setrvačnost
brání předčasné konvergenci
\tb kde je to výhoda: dynamická prostředí, žádná potřeba gradientu, snadná paralelizace}
\most{První provedení principu pracuje ve spojitém prostoru a~pamatuje si trajektorii.}

\blok{2}{Optimalizace rojem částic (PSO)}
\tema{částice $=$ pozice $+$ rychlost $+$ paměť nejlepší navštívené pozice
\tb rovnice rychlosti: setrvačnost $+$ kognitivní složka $+$ sociální složka; aktualizace
pozice
\tb parametry: klesající $w$, $c_1$ a~$c_2$; omezení $v_{\max}$ nebo constriction faktor
\tb topologie gbest vs.\ lbest (kruh, mřížka) a~její vliv na exploraci
\tb rozdíl proti GA: žádné operátory, částice mají paměť, informace teče jednosměrně od
lepších}
\most{Druhé provedení pracuje na grafu a~jako paměť používá samo prostředí.}

\blok{3}{Optimalizace mravenčí kolonií (ACO)}
\tema{řešení se \emph{konstruuje} po krocích; pravděpodobnost přechodu
$p_{ij} \propto \tau_{ij}^\alpha \eta_{ij}^\beta$ (feromon $\times$ heuristika)
\tb aktualizace feromonu: odpařování $+$ uložení úměrné kvalitě
\tb \textbf{AS} (aktualizují všichni) vs.\ \textbf{ACS} (globální update jen nejlepší,
lokální update při průchodu, pseudonáhodné pravidlo) vs.\ \textbf{MAX--MIN} (meze feromonu)
\tb akce démona $=$ lokální prohledávání nad zkonstruovaným řešením
\tb typické úlohy: TSP, směrování, rozvrhování; přirozená vhodnost pro měnící se graf}
\most{Ostatní rojové algoritmy jsou variace na stejné téma s~jinou metaforou.}

\blok{4}{Další rojové algoritmy}
\tema{ABC: dělnice, pozorovatelky, průzkumnice; limit jako restart vyčerpaného zdroje
\tb firefly: přitažlivost klesající se vzdáleností
\tb poznámka ke~kritice \uv{metaforických} algoritmů --- řada novějších je přeznačkovaná
varianta téhož}
\most{Poslední blok je to, na co se komise doptá skoro jistě.}

\blok{5}{Srovnání s~evolučními algoritmy a~kdy co použít}
\tema{co mají společné: populace, stochastika, žádný gradient, snadná paralelizace
\tb v~čem se liší: dědičnost a~operátory (EA) vs.\ pohyb a~paměť částice (PSO) vs.\ postupná
konstrukce a~sdílená paměť v~prostředí (ACO)
\tb volba podle úlohy: spojitá optimalizace $\to$ PSO, DE, CMA-ES; konstrukce na grafu
$\to$ ACO; strukturované genotypy $\to$ GA a~GP
\tb hybridizace s~lokálním prohledáváním je téměř vždy výhodná (vazba na ot.~5)}

%---------------------------------------------------------------------
\section{Lokální prohledávání a~memetické algoritmy}
\zneni{Memetické algoritmy, hill climbing, simulované žíhání.}

\begin{nitbox}[Červená nit]
Znění jmenuje dvě jednoduché lokální metody a~jednu hybridní --- postavit to jako
\textbf{stupňování}: čisté lokální prohledávání a~jeho jediná slabina (uvázne v~lokálním
optimu) $\to$ dva různé způsoby, jak ji obejít (řízená náhoda, paměť) $\to$ memetický
algoritmus jako spojení lokálního prohledávání s~populací. Každá metoda je odpovědí na
konkrétní selhání té předchozí.
\end{nitbox}

\blok{1}{Lokální prohledávání a~horolezecký algoritmus}
\tema{okolí, soused, krajina; co je lokální optimum
\tb varianty: steepest ascent, first improvement, stochastický, s~náhodnými restarty
\tb typické patologie: lokální optima, plošiny, hřebeny}
\most{První lék je dovolit zhoršující krok --- ale řízeně a~stále méně často.}

\blok{2}{Simulované žíhání}
\tema{fyzikální analogie; Metropolisovo kritérium $P = \min(1, e^{-\Delta/T})$
\tb rozvrh chlazení: geometrický v~praxi, logaritmický pro důkaz konvergence
\tb pro pevné $T$ konverguje k~Boltzmannovu rozdělení
\tb krajní případy: $T \to 0$ dá hill climbing, $T \to \infty$ náhodnou procházku}
\most{Druhý lék je opačný: nedovolit vrátit se tam, odkud jsem přišel.}

\blok{3}{Tabu prohledávání a~scatter search}
\tema{tabu list \emph{tahů}, tenure, aspirační kritérium
\tb střednědobá a~dlouhodobá (frekvenční) paměť: intenzifikace vs.\ diverzifikace
\tb vždy se jde do nejlepšího nezakázaného souseda, i~když je horší
\tb scatter search: referenční množina kombinující kvalitu a~diverzitu, aritmetické
kombinace $+$ vylepšení}
\most{Všechny tyhle metody pracují s~jedním řešením --- populace přinese to, co jim chybí.}

\blok{4}{Memetické algoritmy}
\tema{EA $+$ lokální prohledávání nad jedinci: globální explorace $+$ lokální doladění
\tb \textbf{lamarckismus} (vylepšení se zapíše do genotypu, rychlejší konvergence) vs.\
\textbf{baldwinismus} (mění jen fitness, zachovává diverzitu)
\tb návrhová rozhodnutí: na koho, jak často a~jak hluboko lokální prohledávání pouštět
\tb memetic computing: memy v~genomu nebo v~koevolvující populaci}
\most{Poslední věc je společná všem metodám v~téhle otázce: jak nastavit parametry.}

\blok{5}{Ladění a~řízení parametrů}
\tema{ladění (před během) vs.\ řízení (za běhu)
\tb pevné, adaptivní a~samoadaptivní strategie
\tb vazba na pravidlo $1/5$ a~samoadaptaci u~ES (ot.~3)}

%---------------------------------------------------------------------
\section{Aplikace evolučních algoritmů}
\zneni{Aplikace evolučních algoritmů (evoluce expertních systémů, neuroevoluce, řešení
kombinatorických úloh, vícekriteriální optimalizace).}

\begin{nitbox}[Červená nit]
Znění vyjmenovává čtyři aplikace, takže osnova je hotová --- pointa je ve \emph{způsobu}
výkladu. U~každé aplikace odpovědět na \textbf{stejné tři otázky: co je jedinec, co je
fitness, co je na té úloze těžké}. Ta trojice drží celou odpověď a~zároveň ukazuje, že
aplikace nejsou výčet, ale čtyři různé odpovědi na otázku \uv{jak úlohu zakódovat}.
\end{nitbox}

\blok{1}{Evoluce pravidlových a~expertních systémů}
\tema{\textbf{michiganský} přístup (jedinec $=$ jedno pravidlo, řešením je celá populace,
nutné rozdělení odměny) vs.\ \textbf{pittsburghský} (jedinec $=$ celá sada pravidel, fitness
přímočará, operátory složité)
\tb Hollandův LCS: podmínky nad $\{0,1,*\}$, vnitřní zprávy, bucket brigade
\tb ZCS: match set $\to$ výběr akce $\to$ action set; posílení, daň, cover operátor
\tb \textbf{XCS}: klasifikátor (podmínka, akce, predikce, chyba, fitness), \textbf{fitness
podle přesnosti predikce}, aktualizace Q-learningem, nikový GA v~action setu $\Rightarrow$
úplná a~minimální mapa vstup--výstup
\tb umět vysvětlit, proč přesnost jako fitness řeší nedostatky \uv{strength-based} systémů
\tb varianty UCS a~XCSF}
\most{Druhá aplikace nahradí pravidla neuronovou sítí --- a~s~ní se změní, co je genotyp.}

\blok{2}{Neuroevoluce}
\tema{co se evolvuje: váhy / topologie / obojí / hyperparametry učení
\tb proč EA místo gradientu: bez derivací, rekurentní sítě, RL bez rozlišitelné odměny,
paralelizace
\tb \textbf{problém konkurujících konvencí} (permutační problém) $\Rightarrow$ křížení sítí
je bez opatření destruktivní
\tb kódování architektury: přímé (matice), gramatické (Kitano), buněčné (Gruau); GNARL
\tb \textbf{NEAT}: inovační čísla (křížení jen shodných genů), speciace se sdílenou fitness
(ochrana nových topologií), complexification od minimální sítě
\tb \textbf{HyperNEAT}: CPPN generuje váhy ze souřadnic neuronů, zachycuje symetrie
a~regularity
\tb NAS: prostor, strategie, odhad výkonu (sdílení vah); DeepNEAT a~CoDeepNEAT; ES pro deep
RL}
\most{Třetí skupina úloh má diskrétní strukturu a~tvrdá omezení --- tam rozhoduje kódování.}

\blok{3}{Kombinatorické úlohy a~práce s~omezeními}
\tema{batoh: bitová mapa a~čtyři strategie pro omezení --- \emph{penalizace} (volba váhy),
\emph{oprava} (lamarckovsky nebo baldwinovsky), \emph{dekodér} (každý genotyp je přípustný),
\emph{vícekriteriální} formulace
\tb TSP: fitness triviální, problém je kódování; permutační reprezentace $+$ PMX/OX/CX/ER,
inicializace heuristikou, mutace inverzí (2-opt); které reprezentace nepoužívat
\tb SAT: fitness $=$ počet splněných klauzulí (samotné $f$ nestačí), vážení klauzulí,
zjemňující funkce; WalkSAT jako hybridizace
\tb rozvrhování: tvrdá omezení řešit v~operátorech, měkká ve fitness; job shop --- permutace
s~dekodérem vs.\ přímé kódování
\tb obecné pravidlo: rozhoduje kódování a~operátory respektující strukturu úlohy}
\most{Poslední aplikace mění samotný pojem \uv{nejlepší řešení}.}

\blok{4}{Vícekriteriální optimalizace}
\tema{dominance, Pareto množina a~Pareto fronta; dvojí cíl $=$ konvergence $+$ diverzita
\tb skalarizace: lineární (nenajde nekonvexní části fronty), $\varepsilon$-constraint,
Čebyševova
\tb \textbf{NSGA-II}: nedominované třídění do front, crowding distance, crowded comparison,
elitismus spojením rodičů a~potomků
\tb SPEA2 (archiv, síla dominujících, hustota), MOEA/D (skalarizace $+$ váhové vektory
a~sousedství), SMS-EMOA (hypervolume), NSGA-III (referenční body)
\tb metriky: hypervolume (monotónní vůči dominanci, nepotřebuje skutečnou frontu), GD, IGD,
spread
\tb kde se to hodí i~jinde: bloat u~GP jako druhé kritérium (vazba na ot.~1)}

%=====================================================================
\okruh{03}{Dobývání znalostí}{Sedm velkých otázek podle oficiálního znění okruhu
  --- kostry odpovědí}

\section*{Přehled: kostry odpovědí}
\addcontentsline{toc}{section}{Přehled: kostry odpovědí}

\begin{center}
\small
\begin{tabular}{@{}>{\raggedright\arraybackslash}p{3.0cm}>{\raggedright\arraybackslash}p{13.4cm}@{}}
\toprule
\textbf{Otázka} & \textbf{Kostra (bloky za sebou)} \\
\midrule
1 Paradigmata & vymezení a~kořeny $\to$ proces KDD $\to$ typy úloh $\to$ data a~škály
  $\to$ metodologie \\
2 Příprava dat & proč 60--80~\% $\to$ čištění $\to$ integrace a~transformace $\to$
  diskretizace $\to$ redukce $\to$ relevance atributů \\
3 Metody & mapa tří rodin $\to$ asociační pravidla $\to$ učení s~učitelem $\to$ klastrová
  analýza \\
4 Pravidla & proč pravidla $\to$ charakteristická vs.\ diskriminační $\to$ jak je získat
  $\to$ objektivní zajímavost $\to$ subjektivní zajímavost $\to$ aplikace \\
5 Reprezentace & formy reprezentace $\to$ kritéria kvality $\to$ odhad chyby $\to$ míry
  a~ROC $\to$ tři roviny vyhodnocování $\to$ vizualizace \\
6 Sociální sítě & vymezení $\to$ modely sítí $\to$ centralita $\to$ detekce komunit \\
7 Praxe & metody dobývání znalostí (metoda $\to$ oblast) $\to$ metody analýzy sítí
  (metoda $\to$ oblast) \\
\bottomrule
\end{tabular}
\end{center}

%---------------------------------------------------------------------
\section{Základní paradigmata dobývání znalostí}
\zneni{Základní paradigmata dobývání znalostí.}

\begin{nitbox}[Červená nit]
Dobývání znalostí není algoritmus, ale \textbf{proces}: paradigma určuje typ úlohy, typ
úlohy určuje metodu \emph{i}~kritérium, podle kterého se výsledek hodnotí. Celá odpověď je
jedna kauzální linka \emph{proces $\to$ typ úlohy $\to$ data $\to$ metodologie}.
\end{nitbox}

\blok{1}{Vymezení oboru}
\tema{definice KDD (netriviální, implicitní, dosud neznámé, potenciálně užitečné)
\tb data --- informace --- znalost
\tb proč obor vznikl: objem dat vs.\ potřeba podpory rozhodování
\tb tři kořeny: strojové učení, databáze, statistika
\tb dotazování a~OLAP vs.\ KDD}
\most{OLAP odpovídá na otázku, kterou umím položit --- na vzor, o~kterém nevím, potřebuju
proces.}

\blok{2}{Proces KDD}
\tema{pět fází: výběr $\to$ předzpracování $\to$ transformace $\to$ dobývání $\to$
interpretace
\tb dobývání jako \emph{jediný} krok procesu
\tb interaktivnost a~iterativnost
\tb 60--80~\% času na přípravu dat}
\most{Proces je rámec; obsah mu dává to, co vlastně chci najít.}

\blok{3}{Typy úloh --- vlastní paradigmata}
\tema{klasifikace a~predikce \tb popis \tb hledání nugetů
\tb deskriptivní vs.\ prediktivní úlohy
\tb s~učitelem / bez učitele / poloučené
\tb jak typ úlohy mění kritérium úspěchu: přesnost vs.\ srozumitelnost a~pokrytí vs.\
překvapivost}
\most{Volbu úlohy i~metody přitom předem omezuje to, jaká data mám k~dispozici.}

\blok{4}{Data jako vstup}
\tema{tabulka objekty $\times$ atributy jako cílový formát
\tb škály měření: nominální, ordinální, intervalová, poměrová
\tb co která škála dovoluje počítat
\tb strukturovaná vs.\ nestrukturovaná data (text, sekvence, graf, prostor a~čas)}
\most{Aby to byl inženýrský obor a~ne řemeslo, potřebuje proces standardizovanou podobu.}

\blok{5}{Metodologie}
\tema{SEMMA \tb CRISP-DM: šest fází, cyklické, porozumění problému na začátku i~na konci
\tb ASUM-DM (CRISP-DM $+$ projektové řízení)
\tb proč pořadí fází není striktní}

%---------------------------------------------------------------------
\section{Příprava dat, výběr atributů a~analýza jejich relevance}
\zneni{Příprava dat; výběr atributů a~metody pro analýzu jejich relevance.}

\begin{nitbox}[Červená nit]
Cílem přípravy je \textbf{jedna tabulka objekty $\times$ atributy, na které zvolená metoda
funguje}. Každý krok je rozhodnutí, které se dá obhájit i~pokazit --- a~správnost volby se
vždy poměřuje tím, jakou metodu chci použít potom.
\end{nitbox}

\blok{1}{Proč je příprava klíčová}
\tema{60--80~\% času \tb jak reálná data vypadají: chybějící hodnoty, šum, nekonzistence,
více zdrojů, různé škály
\tb příprava je závislá na metodě (rozhodovací strom vs.\ metoda pracující se vzdáleností)}
\most{Nejdřív odstranit vady dat, teprve pak je tvarovat pro metodu.}

\blok{2}{Čištění dat}
\tema{chybějící hodnoty: ignorovat objekt / nová hodnota \uv{nevím} / existující hodnota
(modus, poměrně, náhodně) / doplnit modelem
\tb kdy je absence sama informací
\tb šum a~odlehlé hodnoty \tb kvartily, IQR a~pravidlo $1{,}5\,\mathrm{IQR}$, box plot
\tb odlehlá hodnota vs.\ chyba měření}
\most{Vyčištěná data ještě nejsou srovnatelná --- atributy mají různé jednotky a~rozsahy.}

\blok{3}{Integrace a~transformace}
\tema{spojení zdrojů, duplicity a~redundance
\tb standardizace: desítkové škálování, min--max, z-skóre a~jeho robustní varianta
\tb motivace: bez ní vzdálenost ovládne atribut s~největším rozsahem
\tb derivované a~agregované atributy \tb kódování nominálních atributů}
\most{Některé metody navíc neumí pracovat se spojitými hodnotami vůbec.}

\blok{4}{Diskretizace}
\tema{ekvidistantní \tb ekvifrekvenční
\tb s~využitím třídy: entropie shora dělením, $\chi^2$ zdola slučováním
\tb fuzzy diskretizace: rozostřené hranice, součet příslušností 1, váha objektu
\tb seskupování hodnot kategoriálního atributu (KEX: Assign a~Group)}
\most{Vedle typu hodnot je druhým problémem množství --- objektů i~atributů.}

\blok{5}{Redukce dat}
\tema{příliš mnoho objektů: vzorkování a~reprezentativnost, více modelů nad podmnožinami
\tb nevyvážené třídy: váhy chyb nebo různé pravděpodobnosti výběru
\tb příliš mnoho atributů: \emph{transformace} (PCA --- ztráta interpretace, nutnost znát
všechny původní) vs.\ \emph{selekce} (podmnožina, interpretace zůstává)
\tb prokletí dimenzionality}
\most{Selekce ale předpokládá, že umím říct, který atribut je relevantní --- a~to je
samostatná úloha.}

\blok{6}{Analýza relevance atributů}
\tema{filtry vs.\ wrappery: nezávislost na modelu a~rychlost vs.\ zachycení interakcí a~cena
\tb prohledávání bottom-up (přidávání) a~top-down (odebírání)
\tb kritéria na kontingenční tabulce: $\chi^2$, stupně volnosti, podmínka očekávaných
četností, Fisherův exaktní test pro malé četnosti
\tb entropická kritéria: $H(A)$, vzájemná informace, informační závislost a~proč se
normalizuje
\tb vztahy mezi atributy: korelace a~regrese}

%---------------------------------------------------------------------
\section{Metody pro dobývání znalostí}
\zneni{Metody pro dobývání znalostí; asociační pravidla, přístupy založené na principu
učení s~učitelem a~klastrová analýza.}

\begin{nitbox}[Červená nit]
Tři rodiny, které okruh jmenuje, odpovídají třem otázkám nad daty: \emph{co jde s~čím},
\emph{kam to patří}, \emph{co je si podobné}. U~každé rodiny stačí \textbf{jeden algoritmus
do hloubky} a~ostatní jmenovat jako varianty, které řeší jeho konkrétní slabinu. Je to
nejobsáhlejší otázka okruhu --- osnovu ohlásit dopředu a~nechat komisi vybrat hloubku.
\end{nitbox}

\blok{1}{Mapa (anons)}
\tema{asociační pravidla $=$ deskriptivní, bez učitele
\tb klasifikace $=$ prediktivní, s~učitelem
\tb klastrování $=$ deskriptivní, bez učitele
\tb u~každé rodiny stejný skelet: formulace úlohy $\to$ základní algoritmus $\to$ co ho
tlačí dál}
\most{Začnu od úlohy, která nepotřebuje žádné značky ani cílovou proměnnou.}

\blok{2}{Asociační pravidla a~frekventované vzory}
\tema{analýza nákupního košíku jako motivace
\tb support, confidence, lift (a~co znamená lift $<1$)
\tb velikost prostoru pravidel $O(2^m)$
\tb \textbf{Apriori}: uzavřenost dolů, generování kandidátů a~prořezávání, prohledávání do
šířky
\tb slabiny $\Rightarrow$ \textbf{FP-Growth} (dva průchody, sdílení prefixů, rozděl a~panuj)
\tb problém vzácných položek $\Rightarrow$ \textbf{MS-Apriori} (vlastní MIS, ztráta
uzavřenosti dolů)
\tb varianty úlohy: taxonomie a~virtuální položky, třídní pravidla (CAR), sekvenční vzory
(GSP)
\tb nadměrné množství pravidel jako hlavní praktický problém}
\most{Nadměrné množství pravidel vede na měření zajímavosti (ot.~4); tady pokračuju
k~úlohám, kde cílovou třídu naopak znám.}

\blok{3}{Přístupy založené na principu učení s~učitelem}
\tema{formulace: trénovací množina, generalizace, chyba na nových datech, přeučení
\tb \textbf{rozhodovací stromy} jako reprezentant: TDIDT, volba atributu podle entropie
a~informačního zisku, převod na pravidla, prořezávání
\tb linie ID3 $\to$ C4.5 (poměrný zisk, spojité a~chybějící hodnoty) $\to$ CART (binární
stromy, Gini)
\tb \textbf{bayesovská klasifikace}: MAP, naivní předpoklad podmíněné nezávislosti, nulové
pravděpodobnosti
\tb \textbf{SVM}: maximální okraj, podpůrné vektory, přechod k~duálu, kernelový trik
\tb \textbf{logistická regrese} a~\textbf{diskriminační analýza} (LDA vs.\ QDA) jako
statistický pól
\tb \textbf{neuronové sítě} a~ELM
\tb \textbf{ensembly}: bagging snižuje rozptyl, náhodné lesy odstraňují korelaci komponent,
boosting snižuje zkreslení
\tb srovnávací osy: přesnost, srozumitelnost, cena učení, typ dat, počet tříd}
\most{Když značky nemám vůbec, mizí i~kritérium správnosti --- a~to mění celou úlohu.}

\blok{4}{Klastrová analýza}
\tema{formulace úlohy a~subjektivita pojmu \uv{dobrý klastr}
\tb míry vzdálenosti: $L_1$, $L_2$, $L_\infty$, Minkowského, Mahalanobisova
\tb nutnost normalizace před shlukováním
\tb \textbf{$k$-means}: iterace, monotonie ceny, závislost na inicializaci $\Rightarrow$
$k$-means++
\tb reprezentanti z~dat: $k$-medians, $k$-medoids, PAM $\Rightarrow$ CLARA a~CLARANS
(škálovatelnost)
\tb \textbf{hierarchické shlukování}: dendrogram, typy linkage, volba počtu klastrů řezem;
CURE
\tb \textbf{hustotní a~mřížkové metody}: DBSCAN (libovolný tvar klastru, šum), MAFIA
\tb LVQ jako shlukování s~učitelem
\tb jak se volí $k$}

\lbl{Rezerva:} FP-strom na malém příkladu \tb srovnání náhodných lesů a~boostingu při stejném
počtu stromů \tb složitost indukce a~klasifikace u~stromů \tb spektrální shlukování
(a~upozornit, že se vrátí v~ot.~6).

%---------------------------------------------------------------------
\section{Charakteristická a~diskriminační pravidla a~jejich zajímavost}
\zneni{Metody pro extrakci charakteristických diskriminačních pravidel a~měření jejich
zajímavosti.}

\begin{nitbox}[Červená nit]
Pravidlo je \textbf{nejsrozumitelnější forma znalosti}, proto mu okruh věnuje samostatnou
otázku. Osnova má dvě poloviny: nejprve \emph{dvě zrcadlové role} pravidla (popsat třídu
vs.\ odlišit ji) a~způsoby, jak pravidla získat; pak otázka, \emph{které z~tisíců
vygenerovaných pravidel jsou zajímavá} --- a~že na ni existuje objektivní i~subjektivní
odpověď.
\end{nitbox}

\blok{1}{Proč vůbec pravidla}
\tema{klasifikátor jako matematická funkce vs.\ pravidlo jako výrok
\tb srozumitelnost jako samostatné kritérium kvality
\tb pravidlo v~logice: implikace a~ekvivalence, kdy je booleovská implikace pravdivá
\tb fuzzy implikace jako reziduum t-normy --- \textbf{stupeň pravdivosti, ne
pravděpodobnost}}
\most{Jedno a~totéž pravidlo přitom může sloužit dvěma opačným účelům.}

\blok{2}{Dvě role: charakterizace a~diskriminace}
\tema{charakteristické pravidlo: popisuje cílovou třídu, nutná podmínka, \textbf{t-weight}
\tb diskriminační pravidlo: odlišuje cílovou třídu od kontrastních, postačující podmínka,
\textbf{d-weight}
\tb výpočet obou z~tabulky
\tb d-weight vždy srovnat s~apriorním podílem třídy}
\most{Zbývá, odkud pravidla vzít --- a~cest je několik, každá z~jiné tradice.}

\blok{3}{Jak pravidla získat}
\tema{extrakce z~klasifikačního stromu: cesta $=$ konjunkce, listy téže třídy $=$ disjunkce,
prořezávání pravidel, šikmá dělení a~cena za ně
\tb třídní asociační pravidla (CAR) jako druhá cesta
\tb \textbf{evoluční konstrukce}: fitness $=$ míra kvality pravidla; michiganský přístup
(jedinec $=$ pravidlo) vs.\ pittsburghský (jedinec $=$ soubor pravidel)
\tb \textbf{observační pravidla} a~GUHA: čtyřpolní tabulka, zobecněný kvantifikátor, dvě
rodiny --- odhady pravděpodobností (podpora, konfidence) a~testy hypotéz (Fisherův test)
\tb statistický pohled: regresní analýza jako charakteristický popis vztahu, diskriminační
analýza jako klasifikace prokládáním rozdělení (LDA vs.\ QDA)}
\most{Všechny cesty mají stejný výstup: příliš mnoho pravidel. Odtud druhá polovina otázky.}

\blok{4}{Zajímavost objektivně}
\tema{pravidlo jako binární klasifikátor $\Rightarrow$ přesnost, precision ($=$ konfidence),
recall, TNr, F-míra, AUC
\tb prahování podporou a~konfidencí jako zobecněný kvantifikátor
\tb různé ceny chyb a~střední hodnota ceny
\tb filtrace redundantních a~triviálních pravidel}
\most{Míra ale neřekne, jestli je pravidlo užitečné pro toho, kdo se podle něj rozhoduje.}

\blok{5}{Zajímavost subjektivně}
\tema{srozumitelnost \tb jednoduchost \tb neočekávanost \tb akceschopnost
\tb čtyři role měr v~praxi: fitness, řazení, filtr, práh}
\most{Že to celé funguje, je nejlépe vidět na doménách, kde jsou pravidla nasazená.}

\blok{6}{Aplikace}
\tema{spam (metainformace, karanténa jako třetí třída, ochrana proti falešně pozitivním)
\tb malware (dvouvrstvá hierarchie pravidel)
\tb signatury síťových hrozeb \tb doporučování z~clickstreamu}

\lbl{Rezerva:} Łukasiewiczova disjunkce u~spamu \tb konkrétní čtyřpolní tabulka $(a,b,c,d)$
\tb počty parametrů LDA vs.\ QDA.

%---------------------------------------------------------------------
\section{Reprezentace, vyhodnocování a~vizualizace znalostí}
\zneni{Reprezentace, vyhodnocování a~vizualizace získaných znalostí.}

\begin{nitbox}[Červená nit]
Tři otázky za sebou: \emph{v~čem znalost zapsat}, \emph{jak poznat, že je dobrá}, \emph{jak
ji ukázat člověku}. Napříč všemi třemi jde jedna osa --- \textbf{směna mezi přesností
a~srozumitelností} --- a~pointa je, že vyhodnocení uzavírá cyklus KDD zpět k~datům.
\end{nitbox}

\blok{1}{Formy reprezentace znalostí}
\tema{IF-THEN pravidla (nejsrozumitelnější, lokální) \tb rozhodovací strom (globální,
čitelný do malé hloubky) \tb rozhodovací tabulka \tb matematická funkce (kompaktní,
neinterpretovatelná) \tb prototypy a~centroidy \tb pravděpodobnostní model \tb grafy a~sítě
\tb instance ($k$-NN, CBR)
\tb ke~každé formě: z~jaké metody vzniká a~čím se za srozumitelnost platí}
\most{Formu mám --- ale \uv{dobrý model} zatím není definovaný.}

\blok{2}{Kritéria hodnocení metod}
\tema{prediktivní přesnost \tb efektivita učení i~použití \tb robustnost \tb škálovatelnost
\tb interpretovatelnost \tb kompaktnost modelu
\tb že \uv{lepší} neznamená automaticky \uv{přesnější}}
\most{Přesnost je jediné z~kritérií, které se dá měřit číslem --- a~i~to jen nepřímo.}

\blok{3}{Odhad chyby}
\tema{trénovací vs.\ testovací chyba, přeučení
\tb holdout a~zákaz míchání množin
\tb $k$-násobná křížová validace (obvykle 5 a~10; průměr a~interval spolehlivosti)
\tb leave-one-out pro velmi malá data
\tb validační množina na ladění parametrů algoritmu
\tb stratifikace a~opakování}
\most{Číslo z~validace ale samo neříká, jaké chyby model dělá.}

\blok{4}{Míry a~křivky}
\tema{matice záměn: TP, FN, FP, TN
\tb accuracy a~proč nestačí u~nevyvážených dat (uvést příkladem)
\tb precision, recall, $F_1$ jako harmonický průměr
\tb \textbf{ROC}: TPR proti FPR, konstrukce ze seřazeného žebříčku, diagonála $=$ náhoda
\tb \textbf{AUC} a~důvod její existence: křivky dvou klasifikátorů se mohou křížit
\tb skórování, řazení do přihrádek a~\textbf{lift} $+$ ekonomická úvaha
\tb nákladově citlivé a~vícetřídní vyhodnocení}
\most{Statistická kvalita je jen jedna ze tří rovin, ve kterých se slovo
\uv{vyhodnocování} používá.}

\blok{5}{Tři roviny vyhodnocování}
\tema{statistická kvalita modelu \tb zajímavost vzorů (vazba na ot.~4) \tb hodnocení
uživatelem
\tb čtyři kritéria zajímavosti znalosti: srozumitelnost, platnost, užitečnost, novost
\tb první dvě roviny objektivní, třetí subjektivní}
\most{Uživatelská rovina je právě ta, kde rozhoduje způsob prezentace.}

\blok{6}{Vizualizace}
\tema{tři role v~KDD: explorace (Anscombeho kvartet), vizualizace modelu, prezentace
výsledku
\tb technika podle úlohy: histogram, box plot, scatter a~SPLOM, paralelní souřadnice,
heatmapa, dendrogram, projekce, ROC a~lift
\tb zásady a~zkreslení: osy, barva podle typu dat, datová vs.\ dekorativní složka
\tb omezení projekcí: t-SNE a~UMAP zkreslují globální vzdálenosti}

\lbl{Rezerva:} interval spolehlivosti u~křížové validace \tb odhad přesnosti dolní mezí
u~C4.5 \tb konkrétní výpočet liftu (kolik zásilek se vyplatí poslat).

%---------------------------------------------------------------------
\section{Modely pro analýzu sociálních sítí, centralita a~detekce komunit}
\zneni{Modely pro analýzu sociálních sítí; míry centrality, detekce komunit.}

\begin{nitbox}[Červená nit]
Znění okruhu jmenuje \textbf{tři věci: modely, míry centrality, detekci komunit} --- a~jen
na nich odpověď stavět. Rámec (co je síť a~jak se reprezentuje) zabere pár vět, modely
připraví půdu a~\textbf{dvě jmenované úlohy dostanou hloubku}. Zbytek látky SNA se do
patnácti minut nevejde: zmínit jednou větou, že tam je, a~nechat na doplňující dotaz.
\end{nitbox}

\blok{1}{Vymezení a~reprezentace (krátce)}
\tema{SNA studuje \textbf{vztahy, ne entity a~jejich atributy}
\tb reprezentace: seznam hran, matice sousednosti (u~orientovaných nesymetrická)
\tb typy sítí: orientované, vážené, bipartitní, znaménkové
\tb ego síť vs.\ \uv{celá} síť a~problém specifikace hranic}
\most{Než začnu měřit, musím vědět, jak reálné sítě vypadají --- to jsou ty
\uv{modely} ze znění okruhu.}

\blok{2}{Modely sítí}
\tema{lokální struktura: silné a~slabé vazby, triadický uzávěr, homofilie, most a~zkratkový
most
\tb koheze: hustota, klastrovací koeficient, průměr sítě a~průměrná vzdálenost
\tb malý svět $=$ vysoké klastrování $+$ krátké cesty
\tb šest vlastností reálných komplexních sítí
\tb \textbf{bezškálové sítě}: růst $+$ preferenční připojování, mocninné rozdělení stupňů
\tb důsledek: odolnost proti náhodné poruše vs.\ zranitelnost cíleným útokem na huby;
jádro--periferie}
\most{V~takové síti nejsou uzly rovnocenné --- a~otázka \uv{kdo je důležitý} má víc
odpovědí.}

\blok{3}{Míry centrality}
\tema{stupňová (kolik lidí dosáhnu přímo) \tb mezilehlá (kde se síť rozpadne) \tb blízkostní
(rychlost šíření) \tb vlastní (jak dobře jsem spojen s~dobře spojenými)
\tb na jakou otázku každá odpovídá $\Rightarrow$ proč nestačí jedna míra
\tb vlastní centralita $\to$ PageRank a~HITS jednou větou jako její pokračování
\tb strukturní ekvivalence a~pozice v~síti}
\most{Centralita hledá jednotlivce; druhá jmenovaná úloha hledá skupiny.}

\blok{4}{Detekce komunit (hlavní blok)}
\tema{definice komunity je subjektivní; explicitní vs.\ implicitní skupiny
\tb řezy: minimální a~jeho nevyváženost, poměrový a~normalizovaný řez, NP-těžkost
\tb \textbf{Girvanové--Newman} (divizivní, mezilehlost hran) $+$ \textbf{modularita proti
null modelu}
\tb \textbf{Louvain} (dvě fáze, nevýhody) $\to$ \textbf{Leiden}
\tb taxonomie čtyř kategorií: uzlově, skupinově, síťově a~hierarchicky orientované
\tb překrývající se komunity (CPM), $k$-klika / $k$-jádro / $k$-plex
\tb výzvy: neznámý počet komunit, heterogenní velikosti, vyhodnocení bez pravdy}

\lbl{Mimo znění okruhu --- jen jmenovat, rezerva na doplňující dotaz:} modelování vlivu
(LTM vs.\ ICM, maximalizace vlivu) \tb vliv vs.\ korelace a~shuffle test \tb výběr vs.\
sociální vliv, afiliační sítě a~tři uzávěry \tb Schellingův model \tb strukturní balance
a~věta o~balanci \tb HITS a~PageRank podrobně (tlumicí faktor, markovský řetězec) \tb
predikce vazeb (společní sousedé, Jaccard, Adamic--Adar, Katz) \tb spektrální shlukování
a~maticová faktorizace \tb tematické komunity a~bipartitní jádra.

%---------------------------------------------------------------------
\section{Praktické využití technik}
\zneni{Praktické využití technik pro dobývání znalostí a~analýzu sociálních sítí.}

\begin{nitbox}[Červená nit]
Odpověď je \textbf{procházka metodami z~obou přednášek, kde ke~každé patří konkrétní oblast
využití a~jedna dvě věty o~tom, proč se tam hodí právě ona}. Držet pořadí podle metod
(ne podle domén) --- tím se ukáže, že látka drží pohromadě, a~zároveň se pokryje celý okruh.
\end{nitbox}

\blok{1}{Dobývání znalostí}

\apl{Asociační pravidla (Apriori, FP-Growth)}{analýza nákupního košíku}
Najdou položky kupované společně; výstup se používá na rozmístění zboží, cross-selling
a~doporučování \uv{zákazníci koupili také}. Hodnotí se hlavně liftem --- podpora
a~konfidence samy vyrobí spoustu triviálních pravidel o~častých položkách.

\apl{Sekvenční vzory (GSP)}{web usage mining a~clickstream}
Zajímá pořadí, ne jen souvýskyt: typické cesty webem, posloupnosti nákupů, sledy událostí
v~logu. Slouží k~reorganizaci webu a~k~načasování nabídky.

\apl{Rozhodovací stromy}{kreditní skórování}
Volí se ne kvůli přesnosti, ale kvůli srozumitelnosti --- banka musí zamítnutí zdůvodnit
a~cesta stromem \emph{je} to zdůvodnění. Stejný důvod je za jejich užitím v~medicínské
diagnostice.

\apl{Naivní Bayes}{filtrace spamu}
Zvládne tisíce příznaků (slov) a~levné inkrementální doučování z~pošty označené uživatelem.
Předpoklad podmíněné nezávislosti je nereálný, ale klasifikace na něm překvapivě dobře
funguje.

\apl{SVM}{kategorizace textu}
Řídké vysoce dimenzionální dokumentové vektory jsou přesně to, na co je maximalizace okraje
dobrá. Cena: model není interpretovatelný a~z~principu klasifikuje do dvou tříd.

\apl{Logistická regrese}{skórování rizik (úvěrových, medicínských)}
Dává pravděpodobnost, ne jen třídu, a~koeficienty se čtou jako poměry šancí --- proto je
standardem všude, kde se rozhodnutí musí obhájit před regulátorem.

\apl{Diskriminační analýza (LDA, QDA)}{klasifikace na malých datech}
Když se dá předpokládat normální rozdělení tříd, stačí odhadnout střední hodnoty
a~kovarianční matici. LDA se navíc používá jako redukce dimenze pro zobrazení tříd.

\apl{Neuronové sítě a~ELM}{predikce a~rozpoznávání}
Nasazují se tam, kde je vztah silně nelineární a~interpretace se nevyžaduje. ELM je volba
pro situace, kdy se model musí přetrénovávat rychle a~často.

\apl{Ensembly (náhodné lesy, boosting)}{detekce podvodů}
Na tabulkových datech nejsilnější volba, protože kombinují nízké zkreslení stromů
s~potlačením rozptylu. Pojišťovna s~nimi netvrdí \uv{podvod}, ale \emph{zaostřuje pozornost}
vyšetřovatele na nejrizikovější případy.

\apl{$k$-means}{segmentace zákazníků}
Centroid se čte jako \uv{typický zákazník segmentu} a~je podkladem pro odlišnou nabídku.
Nutná normalizace atributů a~volba $k$ --- obojí je rozhodnutí analytika, ne dat.

\apl{Hierarchické shlukování}{taxonomie, bioinformatika}
Nevyžaduje předem zadaný počet skupin --- ten se volí až řezem dendrogramu, což se hodí
při explorativní analýze (genová exprese, členění trhu).

\apl{DBSCAN}{prostorová data a~detekce anomálií}
Najde shluky libovolného tvaru a~co nikam nepatří, označí za šum. Proto se používá na GPS
stopy, mapy incidentů a~hledání odlehlých pozorování.

\apl{LVQ}{rozpoznávání vzorů a~klasifikace signálů}
Reprezentanti tříd se uloží jako pár vektorů a~klasifikace je pak levná --- výhoda
v~zařízeních s~omezenými prostředky.

\apl{Regrese a~posuvné okno}{predikce poptávky a~časové řady}
Predikce se převede na učení s~učitelem tím, že vstupem je okno minulých hodnot. Používá se
na plánování zásob, spotřeby energie a~zatížení systémů.

\apl{PCA a~selekce atributů}{předzpracování a~vizualizace}
Sníží dimenzi před shlukováním nebo klasifikací a~umožní data zobrazit. PCA za to platí
ztrátou interpretace, selekce atributů ji zachová.

\apl{Charakteristická a~diskriminační pravidla}{malware a~signatury síťových hrozeb}
Detekce musí být rychlá a~auditovatelná, proto se znalost udržuje v~podobě pravidel, ne
černé skříňky. Pravidla se dají ručně opravit a~verzovat.

\apl{Skórování, lift a~ROC}{cílený marketing a~churn}
Nejde o~třídu, ale o~pořadí: model seřadí zákazníky podle rizika odchodu a~lift křivka
odpoví na ekonomickou otázku, kolik jich má smysl oslovit. Data jsou nevyvážená, takže
accuracy je bezcenná.

\apl{Detekce odlehlých hodnot}{kontrola kvality a~fraud v~provozu}
Odlehlá hodnota je jednou chyba měření a~podruhé právě to hledané --- rozlišení je věcná
otázka domény, ne statistiky.

\blok{2}{Analýza sociálních sítí}

\apl{Stupňová centralita}{hledání nejviditelnějších aktérů}
Nejlevnější míra: kdo má nejvíc přímých kontaktů; v~online sítích slouží k~výběru
influencerů. Sama o~sobě klame --- vysoký stupeň neznamená strategickou pozici.

\apl{Mezilehlá centralita}{brokeři a~kritická místa infrastruktury}
Uzel na mnoha nejkratších cestách drží síť pohromadě: v~kriminální síti je to spojka mezi
buňkami, v~přenosové soustavě místo, jehož výpadek síť rozpojí.

\apl{Blízkostní centralita}{rychlost šíření a~umístění zdroje}
Kdo je v~průměru všem nejblíž, rozšíří zprávu nejrychleji. Používá se při volbě umístění
skladu, serveru nebo osoby, přes kterou se informace pouští do organizace.

\apl{Vlastní centralita}{prestiž v~citačních a~vlastnických sítích}
Důležitost se dědí od důležitých sousedů, ne z~počtu vazeb --- odtud i~míra vlivu firem
propojených vlastnickými podíly.

\apl{PageRank}{řazení výsledků vyhledávání}
Dotazově nezávislá globální míra počítaná offline, proto je v~době dotazu zadarmo a~odolná
proti spamu (vstupní odkazy se nedají snadno vyrobit). Používá se i~mimo web ---
na doporučování a~hodnocení publikací.

\apl{HITS}{tematické vyhledávání --- huby a~autority}
Rozliší rozcestníky od zdrojů, což čistý PageRank nedělá. Počítá se až nad výsledkem
dotazu, a~to je zároveň jeho hlavní provozní slabina.

\apl{Detekce komunit (Louvain)}{segmentace zákazníků a~doporučování}
Na sítích s~miliony uzlů najde skupiny, které se v~atributech zákazníků vůbec neprojeví
(např. jazykové komunity u~operátora). Komunita pak nahradí ručně definovaný segment.

\apl{Modularita a~Girvanové--Newman}{objevování skupin v~menších sítích}
Modularita zároveň odpovídá na otázku, kolik komunit vlastně brát --- proto slouží
i~jako kritérium při vyhodnocování jiných metod.

\apl{Překrývající se komunity (CPM)}{role jednoho člověka ve~více skupinách}
Člověk patří do rodiny, práce i~zájmové skupiny současně; tvrdé dělení tuhle informaci
zahodí, což vadí právě u~doporučování obsahu.

\apl{Predikce vazeb (společní sousedé, Adamic--Adar, Katz)}{\uv{lidé, které možná znáte}}
Stejná úloha řeší i~doporučování zboží (bipartitní síť zákazník--produkt) a~predikci budoucí
spolupráce ve vědeckých sítích. V~řídkých sítích se místo sousedských měr sahá po Katzovi.

\apl{Modelování vlivu (LTM, ICM) a~maximalizace vlivu}{virální marketing}
Otázka zní, kterým $k$~lidem dát produkt zdarma, aby se rozšířil co nejdál. Úloha je
NP-těžká, ale hladový výběr má garantovanou kvalitu, takže je prakticky použitelný.

\apl{Bezškálovost a~odolnost sítě}{cílená vakcinace a~ochrana infrastruktury}
Náhodné výpadky síť snese, cílený zásah do hubů ji rozloží. Odtud plyne očkování podle počtu
kontaktů místo plošného i~to, proč se internetové viry nedají vymýtit.

\apl{Malý svět a~modely šíření}{epidemiologie}
Krátké cesty vysvětlují, proč se nákaza --- a~stejně tak informace nebo panika --- šíří
mnohem rychleji, než by odpovídalo geografické vzdálenosti.

\apl{Znaménkové sítě a~strukturní balance}{polarizace a~predikce znaménka vztahu}
V~hodnotících komunitách se predikuje důvěra a~nedůvěra, historicky se stejným aparátem
analyzovaly evropské aliance před první světovou válkou.

\apl{Homofilie a~rozlišení výběru od vlivu}{vyhodnocování kampaní a~intervencí}
Bez rozlišení, jestli si podobní lidé jen našli jeden druhého, nebo se navzájem ovlivnili,
se dopad kampaně systematicky nadhodnocuje. Nástroj: shuffle test a~longitudinální data.

\apl{Síla slabých vazeb a~mosty}{difuze informací a~pracovní trh}
Nové informace přicházejí přes slabé vazby, ne z~blízkého okruhu, kde všichni vědí totéž.
Z~toho se odvozuje jak strategie oslovování, tak návrh doporučovacích systémů.

\apl{Kombinace centralit}{klíčoví hráči v~kriminálních a~teroristických sítích}
Vůdce nemusí mít nejvyšší stupeň --- bývá skrytý a~chráněný, zato má vysokou mezilehlost.
Proto se pracuje s~více mírami najednou.

\end{document}
